\section{Conclusion}\label{sec:conclusion}

AI decision models read financial news well enough to forecast how a stock will open, and not well enough to
forecast what it does after that. A fine-tuned open-weights model that looked convincing on two validation
quarters did not survive a blinded test; the commercial models kept a signal, but only in the overnight gap,
where the market had already priced the news before anyone could trade on it at the open, as an efficient
market would \citep{fama1970}. The models could not
tell an overreaction from an underreaction, their calls of underreaction never caught up, and only the largest
moves reversed, slightly. No pre-registered strategy survived costs and selection. Out of sample, news explains
a fraction of a percent of a stock's own daily variance, and the commercial models' stated probabilities overstate
their accuracy.

Two lessons reach beyond these models. First, validation on a few quarters is weak evidence: the strongest
exploratory pattern we found, with in-sample Sharpe ratios up to \res{cb-valb-gross-max}, reversed on the test. A frozen plan, a
trial ledger and a single unblinding are what made that visible. Second, masking is not anonymity: commercial
models recognize most masked firms, so backtests of language models on historical news should measure
recognition directly and, where possible, test on news published after the model's release, as our
pre-registered forward test will. The market reads the news first.

\section*{Acknowledgments and statements}
\placeholder{Acknowledgments, funding and conflicts of interest, to be provided by the authors.}
\placeholder{Statement on the use of AI tools, to be confirmed by the authors: an AI coding and research assistant
implemented and ran the pre-registered pipeline under the authors' direction, and drafted this manuscript from
the verified outputs; the authors are responsible for the content.}
